% =========================================================
% II. RELATED WORK
% =========================================================
\section{Related Work}
\label{sec:related_work}

\subsection{Retrieval-Augmented Generation}

Retrieval-Augmented Generation (RAG) kết hợp mô hình ngôn ngữ với một thành
phần truy xuất bên ngoài nhằm cung cấp bằng chứng liên quan trước khi sinh câu
trả lời \cite{lewis2020retrieval,gao2024ragsurvey}. Các phương pháp truy xuất
dense như DPR \cite{karpukhin2020dense} giúp cải thiện khả năng tìm kiếm tài
liệu theo ngữ nghĩa. Các nghiên cứu tiếp theo đã mở rộng theo hướng tổng hợp
nhiều đoạn truy xuất, truy xuất thích ứng kết hợp tự phản hồi, và tổ chức truy
xuất theo cấu trúc phân cấp
\cite{izacard2021fid,asai2024selfrag,sarthi2024raptor}. Tuy nhiên, phần lớn hệ
thống RAG truyền thống vẫn cung cấp bằng chứng truy xuất chủ yếu dưới dạng ngữ
cảnh văn bản thay vì một cấu trúc suy luận tường minh.

Cách tiếp cận này phù hợp với các câu hỏi có thể trả lời từ một hoặc một số đoạn
văn trực tiếp, nhưng gặp hạn chế trong các bài toán đòi hỏi kết nối nhiều bước
suy luận.

\subsection{Multi-hop Question Answering}

Multi-hop Question Answering yêu cầu hệ thống kết hợp thông tin từ nhiều bằng
chứng để suy ra câu trả lời cuối cùng. Các bộ dữ liệu như HotpotQA
\cite{yang2018hotpotqa} và MuSiQue \cite{trivedi2022musique} cho thấy việc truy
xuất đúng từng tài liệu chưa đủ; hệ thống còn phải xác định được chuỗi bằng
chứng liên kết giữa chúng.

Trong miền pháp luật, yêu cầu này đặc biệt quan trọng vì một kết luận pháp lý
thường phụ thuộc vào nhiều điều kiện và hậu quả trung gian. Do đó, việc biểu
diễn rõ cấu trúc suy luận có thể giúp hệ thống hạn chế các bước nối suy luận
không được hỗ trợ bởi văn bản pháp luật.

\subsection{Graph-based and Causal RAG}

Một hướng nghiên cứu gần đây sử dụng đồ thị để tổ chức tri thức và hỗ trợ truy
xuất nhiều bước. GraphRAG khai thác cấu trúc đồ thị nhằm mở rộng khả năng truy
xuất và tổng hợp thông tin trên nhiều thực thể hoặc đoạn văn liên quan
\cite{edge2024graphrag}. Think-on-Graph tiếp tục cho thấy việc duyệt tường minh
trên đồ thị tri thức có thể hỗ trợ suy luận nhiều bước và tạo ra các đường suy
luận có khả năng truy vết \cite{sun2024thinkongraph}.

Gần với hướng nghiên cứu của chúng tôi hơn, CausalRAG sử dụng quan hệ nhân quả
để hướng dẫn quá trình truy xuất và suy luận thay vì chỉ dựa trên độ tương đồng
ngữ nghĩa \cite{wang2025causalrag}. CC-RAG cũng xây dựng các chuỗi nhân quả
nhằm hỗ trợ suy luận nhiều bước có cấu trúc trên các kho dữ liệu chuyên miền
\cite{parekh2025ccrag}. Những hướng nghiên cứu này cho thấy cấu trúc quan hệ
giữa các sự kiện có thể cung cấp tín hiệu bổ sung cho các câu hỏi nhiều bước.

Tuy nhiên, trong bài toán pháp luật, quan hệ được sử dụng trong nghiên cứu này
được hiểu là quan hệ có hướng giữa điều kiện và hệ quả pháp lý được trích xuất
từ văn bản luật, không phải quan hệ nhân quả được ước lượng từ dữ liệu quan sát.

\subsection{Counterfactual Reasoning and Legal QA}

Counterfactual reasoning đã được sử dụng để đánh giá liệu một kết luận còn giữ
nguyên khi một điều kiện hoặc sự kiện trung gian thay đổi. Trong Multi-hop QA,
các nghiên cứu về câu hỏi phản thực tế cho thấy mô hình có thể đạt kết quả đúng
nhưng vẫn dựa vào chuỗi suy luận không phù hợp \cite{guo2023counterfactual}.

Trong lĩnh vực pháp luật, các benchmark như LexGLUE, CaseHOLD và LegalBench,
cùng với các mô hình chuyên miền như LEGAL-BERT, đã cho thấy những đặc trưng
riêng của ngôn ngữ pháp luật cũng như những thách thức còn tồn tại đối với suy
luận pháp lý có cấu trúc
\cite{chalkidis2022lexglue,zheng2021casehold,guha2023legalbench,chalkidis2020legalbert}.
Điều này tạo động lực cho các phương pháp kết hợp truy xuất bằng chứng với biểu
diễn quan hệ pháp lý và cấu trúc suy luận tường minh hơn.

\subsection{Logic programming ba giá trị và biểu diễn đại số tuyến tính}

Logic ba giá trị phân biệt true, false và trạng thái undefined/unknown
\cite{kleene1952metamathematics}. Ngữ nghĩa Kripke--Kleene của Fitting cung cấp
cách xây dựng điểm bất động cho logic program khi thông tin chưa đầy đủ
\cite{fitting1985kripke}. Các nghiên cứu gần đây mã hóa ngữ nghĩa logic program
hai và ba giá trị trong không gian vector \cite{sato2020threevalued}, đồng thời
dùng ma trận và tensor bậc ba cho definite, disjunctive và normal logic program
\cite{sakama2021tensor}. Các công trình này tạo nền tảng cho formulation đại số
của chúng tôi. Runtime hiện tại vẫn là engine quy tắc thưa ký hiệu: AND nằm
trong một quy tắc và OR nằm giữa các quy tắc độc lập có cùng consequent, thay vì
thực hiện phép nhân tensor.

\subsection{Research Gap}

Các phương pháp RAG truyền thống chủ yếu tối ưu khả năng truy xuất bằng chứng
có liên quan về mặt ngữ nghĩa, trong khi các phương pháp dựa trên đồ thị tập
trung vào việc kết nối nhiều thực thể hoặc sự kiện. Tuy nhiên, vẫn còn ít nghiên
cứu kết hợp đồng thời:

\begin{itemize}
    \item biểu diễn các quy tắc pháp luật dưới dạng quan hệ điều kiện--hệ quả;
    \item truy xuất và xếp hạng các đường suy luận nhiều bước;
    \item xác minh vai trò của các sự kiện trung gian bằng phân tích phản thực tế;
    \item và sử dụng kết quả xác minh để tinh lọc bằng chứng trước khi sinh câu trả lời.
\end{itemize}

Nghiên cứu này hướng tới khoảng trống trên bằng cách kết hợp causal graph
retrieval, multi-hop path reasoning và counterfactual verification trong một
pipeline RAG dành cho hỏi đáp pháp luật nhiều bước.